\documentclass[12pt]{article}
\usepackage[utf8]{inputenc}
\usepackage[T1]{fontenc}
\usepackage[margin=1in]{geometry}
\usepackage{setspace}
\usepackage{xcolor}
\usepackage[colorlinks=true, urlcolor=blue!60!black, linkcolor=blue!60!black]{hyperref}
\setlength{\parindent}{2em}
\setlength{\parskip}{1.25em}
\renewcommand{\baselinestretch}{1.0}
\begin{document}
\begin{center}
{\LARGE \textit{Obergefell v. Hodges}, 576 U.S. 644 (2015)\par}
\vspace{0.8em}
{\large Case Brief --- Four-Page Edition\par}
\end{center}
\vspace{1em}

\section{Memory Jogger}

Rejecting a rigid history-and-tradition ceiling for substantive due process, the Supreme Court held 5--4 that the interlocking guarantees of the Fourteenth Amendment's Due Process and Equal Protection Clauses require every state to license same-sex marriages and recognize same-sex marriages lawfully performed out of state (\href{https://www.courtlistener.com/opinion/2812209/obergefell-v-hodges/}{576 U.S. 644}).

\section{Detailed Case Facts}

Fourteen same-sex couples and two widowers challenged laws in Michigan, Kentucky, Ohio, and Tennessee defining marriage as between one man and one woman and refusing to recognize out-of-state same-sex marriages. James Obergefell and John Arthur, together for over two decades, married on a medical-transport plane's tarmac in Maryland after Arthur's ALS diagnosis; when Arthur died, Ohio barred listing Obergefell as surviving spouse, leaving them ``strangers even in death.'' April DeBoer and Jayne Rowse, Michigan nurses, adopted three special-needs children, but each child could have only one legal mother. Army Reserve Sergeant Ijpe DeKoe married Thomas Kostura in New York before deploying to Afghanistan, and Tennessee refused to recognize the marriage.

\section{Procedural History}

Every district court in six suits ruled for the plaintiffs. A divided Sixth Circuit panel (Judge Sutton, over Judge Daughtrey's dissent) reversed in \textit{DeBoer v. Snyder}, holding itself bound by \textit{Baker v. Nelson} (1972) and leaving marriage to the voters, which split from four other circuits. The Supreme Court granted certiorari on whether states must license and recognize same-sex marriages, and decided the case on June 26, 2015.

\section{Judicial Votes}

Reversed, \textbf{5--4}.
\begin{itemize}
\item \textbf{Majority:} Justice Kennedy, joined by Justices Ginsburg, Breyer, Sotomayor, and Kagan.
\item \textbf{Dissents:} Chief Justice Roberts (joined by Scalia and Thomas); Justice Scalia (joined by Thomas); Justice Thomas (joined by Scalia); Justice Alito (joined by Scalia and Thomas).
\end{itemize}

\section{Holding}

Under the Due Process and Equal Protection Clauses, the right to marry is fundamental and same-sex couples may not be deprived of it. States must license same-sex marriages on the same terms as opposite-sex marriages and recognize those lawfully performed elsewhere; \textit{Baker v. Nelson} is overruled. Courts identify fundamental rights through ``reasoned judgment,'' guided but not bounded by history, rather than \textit{Glucksberg}'s narrow formula. Assurances to religious objectors, the ``two consenting adults'' limit, and the open question of heightened scrutiny for sexual orientation are not holdings.

\section{Analysis of Opinions}

\textbf{Majority (Kennedy).} Marriage has endured while changing, as with the end of coverture. Invoking Harlan's \textit{Poe v. Ullman} dissent, Kennedy rejected framing the claim as a narrow ``right to same-sex marriage'': \textit{Loving}, \textit{Zablocki}, and \textit{Turner} asked about the ``right to marry in its comprehensive sense.'' Four principles apply equally to same-sex couples: marriage is an intimate act of autonomy; it supports a unique two-person union beyond \textit{Lawrence}; it safeguards children and families; and it is a keystone of social order to which states attach a vast ``constellation of benefits.'' The clauses' ``synergy'' meant exclusion both burdened liberty and worked a subordinating inequality. Because ``fundamental rights may not be submitted to vote,'' the Court would not await democratic change.

\textbf{Roberts (dissenting).} The question is ``who decides.'' Invoking \textit{Dred Scott} and \textit{Lochner}, he accused the majority of ``effectively overrul[ing] \textit{Glucksberg}'' to enact its own moral philosophy, found no principled limit against plural marriage, and noted that the majority protected the right to ``advocate'' and ``teach'' but omitted ``exercise.''

\textbf{Scalia (dissenting).} A ``judicial Putsch'' and ``threat to American democracy'': every state in 1868 limited marriage to a man and a woman, and the Court's prose offered ``the mystical aphorisms of the fortune cookie.''

\textbf{Thomas (dissenting).} Substantive due process is a ``dangerous fiction''; liberty means freedom from restraint, not entitlement to government benefits, and dignity ``cannot be taken away by the government.''

\textbf{Alito (dissenting).} Same-sex marriage is not ``deeply rooted in this Nation's history and tradition''; states could rationally keep a procreation-linked model, and believers would be vilified.

\section{Examples: Future Applications}

\begin{enumerate}
\item \textbf{Same-Side (Invalid): Birth Certificates.} A state lists a birth mother's husband but not her wife on a donor-conceived child's certificate. This withholds the marital ``constellation of benefits'' (\textit{Pavan v. Smith}).
\item \textbf{Same-Side (Invalid): Spousal Benefits.} A city denies health and survivor benefits to firefighters' lawfully married same-sex spouses.
\item \textbf{Opposite-Side (Upheld): Assisted-Suicide Ban.} Kennedy distinguished \textit{Glucksberg} as possibly appropriate for that right, and \textit{Dobbs} reaffirmed its test.
\item \textbf{Opposite-Side (Upheld): Age and Consanguinity Limits.} \textit{Obergefell} reached only ``two consenting adults whose marriages would pose no risk of harm to themselves or third parties.''
\item \textbf{Fence-Sitter: Plural Marriage.} Three committed co-parenting adults invoke autonomy and child welfare, as Roberts warned; but Kennedy anchored marriage in a two-person union, and plural marriage raises distinct custody and inheritance complexities.
\end{enumerate}

\section{Critique}

\textbf{Originalist and textualist views.} Echoing the dissents, critics argue that nothing in the 1868 meaning required same-sex marriage and that ``reasoned judgment'' replaces historical constraint with judicial values. After \textit{Dobbs}, Justice Thomas urged reconsidering \textit{Obergefell}, prompting the Respect for Marriage Act of 2022.

\textbf{Progressive and doctrinal views.} Kenji Yoshino called the decision a ``game-changer'' for substantive due process but noted it left no tier of scrutiny for sexual orientation; Laurence Tribe praised its ``equal dignity'' fusion of liberty and equality. Family-law scholars faulted its glorification of marriage, which demeans nonmarital families.

\section{Key Quotations}

\begin{quote}
``The identification and protection of fundamental rights is an enduring part of the judicial duty to interpret the Constitution. That responsibility, however, `has not been reduced to any formula.' \textit{Poe v. Ullman}, 367 U.S. 497, 542 (1961) (Harlan, J., dissenting). Rather, it requires courts to exercise reasoned judgment in identifying interests of the person so fundamental that the State must accord them its respect. [\dots] History and tradition guide and discipline this inquiry but do not set its outer boundaries. [\dots] That method respects our history and learns from it without allowing the past alone to rule the present.''
\end{quote}
\noindent --- Justice Kennedy, Opinion of the Court, 576 U.S. at 664.

\begin{quote}
``No union is more profound than marriage, for it embodies the highest ideals of love, fidelity, devotion, sacrifice, and family. In forming a marital union, two people become something greater than once they were. As some of the petitioners in these cases demonstrate, marriage embodies a love that may endure even past death. It would misunderstand these men and women to say they disrespect the idea of marriage. Their plea is that they do respect it, respect it so deeply that they seek to find its fulfillment for themselves. Their hope is not to be condemned to live in loneliness, excluded from one of civilization's oldest institutions. They ask for equal dignity in the eyes of the law. The Constitution grants them that right.''
\end{quote}
\noindent --- Justice Kennedy, Opinion of the Court, 576 U.S. at 681.

\begin{quote}
``If rights were defined by who exercised them in the past, then received practices could serve as their own continued justification and new groups could not invoke rights once denied.''
\end{quote}
\noindent --- Justice Kennedy, Opinion of the Court, 576 U.S. at 671.

\end{document}
